% Hilbert 空间投影的几何：(a) 闭凸集 C（凸多边形示意）上，z-p 是 p 处的外法方向，与任一 u-p 的夹角不小于直角，虚线为过 p 的支撑线；
% (b) 闭子空间 M 上，残差 z-p 与 M 正交。
\begin{tikzpicture}[note/.style={font=\FigureNoteFont\fontsize{7.5}{9}\selectfont},
  dot/.style={circle,fill=FigureInk,inner sep=0pt,minimum size=3pt},
  vec/.style={line width=.9pt,->,shorten >=1.5pt}]
  \begin{scope}
    \filldraw[fill=FigureBlueFill,draw=FigureBlue,line width=.9pt]
      (0,0)--(2.2,-.3)--(2.8,1.2)--(1.4,2.0)--(-.2,1.2)--cycle;
    \coordinate (p) at (2.5,.45);
    \coordinate (z) at (3.7,-.03);
    \coordinate (u) at (.8,.9);
    \draw[FigureMuted,line width=.5pt,dashed] ($(p)!-1.0!(2.8,1.2)$)--($(p)!2.3!(2.8,1.2)$);
    \draw[FigureYellow,vec] (p)--(z);
    \draw[FigureGreen,vec] (p)--(u);
    \node[dot] at (p) {}; \node[dot] at (z) {}; \node[dot] at (u) {};
    \node[right=1pt] at (z) {$z$};
    \node[left=2pt] at (p) {$p$};
    \node[above left=-1pt] at (u) {$u$};
    \node[FigureBlue] at (1.5,1.45) {$C$};
    \node[note,text=FigureYellow!70!FigureInk] at (3.25,.5) {$z-p$};
    \node[note,text=FigureGreen!80!FigureInk] at (1.5,.45) {$u-p$};
    \node[note] at (1.3,-.85) {(a) 闭凸集：夹角不小于直角};
  \end{scope}
  \begin{scope}[xshift=5.4cm]
    \draw[FigureBlue,line width=.9pt] (-.3,-.1)--(3.2,1.0) node[right]{$M$};
    \coordinate (p2) at (2.0,.6);
    \coordinate (z2) at ($(p2)!1.35cm!90:(3.2,1.0)$);
    \coordinate (u2) at (.5,.1268);
    \draw[FigureYellow,vec] (p2)--(z2);
    \draw[FigureGreen,vec] (p2)--(u2);
    \draw[FigureStroke,line width=.5pt] ($(p2)!.2cm!(z2)$) -- ($($(p2)!.2cm!(z2)$)!.2cm!90:(z2)$) -- ($(p2)!.2cm!(u2)$);
    \node[dot] at (p2) {}; \node[dot] at (z2) {}; \node[dot] at (u2) {};
    \node[above] at (z2) {$z$};
    \node[below right=-1pt] at (p2) {$p$};
    \node[below=1pt] at (u2) {$u$};
    \node[note] at (1.4,-.85) {(b) 闭子空间：残差与 $M$ 正交};
  \end{scope}
\end{tikzpicture}
